\documentclass[conference]{IEEEtran}
\IEEEoverridecommandlockouts
\usepackage{cite}
\usepackage{amsmath,amssymb,amsfonts}
\usepackage{algorithmic}
\usepackage{graphicx}
\usepackage{textcomp}
\usepackage{xcolor}
\usepackage{tikz}
\usetikzlibrary{arrows.meta,positioning,fit,backgrounds,shapes.geometric}
\def\BibTeX{{\rm B\kern-.05em{\sc i\kern-.025em b}\kern-.08em
    T\kern-.1667em\lower.7ex\hbox{E}\kern-.125emX}}
\begin{document}

\title{Quality-Adaptive Loss Reweighting for Robust Skin Lesion Classification}

\author{
\IEEEauthorblockN{1\textsuperscript{st} Given Name Surname}
\IEEEauthorblockA{\textit{Department Name} \\
\textit{Institution Name}\\
City, Country \\
email@example.com}
}

\maketitle

\begin{abstract}
Diagnostic image quality varies widely in clinical practice, yet quality-aware training methods have had to estimate quality from image statistics or feature norms rather than observe it, because expert judgments of image quality are almost never recorded alongside diagnostic labels. This paper has proposed quality-adaptive loss reweighting, which scales each training sample's classification loss by a real expert-assigned rating of its diagnostic image. Two opposing directions have been evaluated: down-weighting low-quality samples as less reliable, and up-weighting them to force the model to learn from difficult images. The second, the opposite of what the closest face recognition precedents apply, has improved balanced accuracy by 0.037 and malignant sensitivity by 0.092 over an identical baseline. Because a gain of this size is comparable to the fold-to-fold standard deviation, we have validated it causally rather than statistically: a shuffled-quality control, preserving the weight distribution but permuting the lesion-to-rating assignment, has removed the gain entirely and returned performance to baseline, establishing that the improvement has depended on the quality signal's information content rather than on generic reweighting. The full pipeline has reached 0.840 balanced accuracy, of which 65\% of the total gain has been attributable to the quality-adaptive loss and the remainder to generic evaluation-time averaging. The method has been made possible by MCR-SL, the only public skin lesion dataset recording per-image expert quality ratings, on which we have also reported the first published benchmark and six robustness analyses, three of which have reversed an apparent positive result.
\end{abstract}

\begin{IEEEkeywords}
skin lesion classification, multi-modal fusion, dermoscopy, image quality, metadata fusion, deep learning, medical imaging benchmark
\end{IEEEkeywords}

\section{Introduction}
Medical images acquired in routine practice vary substantially in diagnostic quality. Lighting, focus, framing, and hair occlusion all degrade the evidence a classifier receives, and training as though every image were equally informative assumes a uniformity clinical data does not have. Face recognition has addressed this by making the objective quality aware: AdaFace \cite{adaface} and MagFace \cite{magface} adapt each sample's margin from an estimate of image quality derived from the feature norm, both reducing emphasis on difficult low-quality samples. That reliance on an estimated proxy is a necessity, not a preference, because face datasets do not record what an expert thought of each photograph.

Skin lesion classification has followed a parallel path. Convolutional networks have reached dermatologist-level accuracy on large single-modality archives \cite{esteva2017,ham10000}, and a smaller line of work has fused images with patient metadata \cite{metablock,paduf20}. In neither case has a per-image expert judgment of quality been available, so quality has remained something to be inferred rather than observed.

MCR-SL \cite{mcrsl} has removed that constraint, and is the reason the method proposed here is possible at all. It has paired 779 clinical and 1352 dermoscopic images across 240 lesions from 60 subjects with twenty-two subject-level metadata fields and, distinctively, with a one-to-ten quality rating that dermatologists have assigned to the specific image used for each lesion's diagnosis. It is, to our knowledge, the only public skin lesion dataset carrying such a field, which makes it the enabling dataset for a loss function driven by real expert quality rather than an estimated proxy, not merely a convenient testbed for a pre-existing method.

This work has posed two separate questions: whether performance is degraded by lower-quality images, and whether the quality signal can be used by the training objective to improve robustness.

\subsection{Research Gap}
No prior work has been found that uses a direct expert rating of image quality as a
loss-reweighting signal in medical image classification. Existing quality-aware objectives
instead rely on estimated proxies for image quality. Quality-aware medical imaging methods
have modified the architecture rather than the loss function, leaving how individual
low-quality samples are weighted during training unaddressed. The closest skin-lesion
loss-weighting scheme derives its difficulty signal from the model itself rather than from a
human judgment of the photograph, capturing samples the model finds hard rather than samples
that are genuinely poor in quality.

\subsection{Contributions}
This paper has made the following contributions.
\begin{itemize}
\item \textbf{Quality-adaptive loss reweighting:} Each sample's classification loss has been scaled by a per-sample weight derived from an expert-assigned image-quality rating, allowing lower-quality samples to be up-weighted rather than down-weighted during training.
\item \textbf{Causal validation via a shuffled-quality control:} The quality ratings have been randomly shuffled across lesions while keeping the same distribution of weights. This has removed the entire performance gain, showing that the improvement has come from the actual quality information and not simply from reweighting samples by some arbitrary number.
\item \textbf{Ablation Study:} A five-fold stratified subject-disjoint evaluation protocol has been established on MCR-SL, along with a sweep of 11 ablation configurations that has identified which interventions help at this scale and which do not.
\end{itemize}
%%%%%%%%%%%%%%%%%%%%%%%%%%%%%%%%%%%%%%%%%%%%%%%%%%%%%%%%%
\section{Literature Review} %% Add metrics %%
\label{sec:lit}
Image-only dermoscopic classification has been studied extensively since Esteva \emph{et al.} \cite{esteva2017} matched dermatologist performance on a curated archive, with HAM10000 \cite{ham10000} becoming a standard benchmark. These studies have treated the image as the sole input.

Metadata fusion is a distinct thread. Pacheco and Krohling \cite{metablock} have proposed the Metadata Processing Block, which reweights image features using patient metadata, and report 0.765 balanced accuracy on the six-class PAD-UFES-20 task \cite{paduf20} (2298 clinical images, 1641 lesions). That dataset is the closest comparator in problem shape and is our primary external reference point, discussed later; on its binary cancer-versus-non-cancer form the strongest verified result is 0.836 \cite{diffmic}.

Quality-aware training has been developed most fully in face recognition. AdaFace \cite{adaface} adapts the angular margin using the feature norm as a quality proxy, reducing emphasis on hard examples when an image is estimated to be low quality, since a degraded and difficult face is likely unidentifiable and forcing separation would fit noise. MagFace \cite{magface} ties feature magnitude to quality similarly. Both depend on an estimated proxy rather than a recorded human judgment.

Within medical imaging, quality has more often been handled architecturally. QGDR \cite{qgdr} routes fundus images through quality-conditioned expert branches for diabetic retinopathy grading, reaching 78.32\% accuracy on EyeQ and 80.85\% on DDR, and reporting preserved performance on the lowest-quality stratum. It modifies the architecture rather than the loss, and derives its quality states largely from predicted pseudo-labels rather than from expert annotation.

Closest to this work, Hu and Yang \cite{hu2023} have applied knowledge-based loss weighting on HAM10000, assigning weights at the class level, favouring minority classes, and at the instance level, favouring challenging samples. Their instance-level term shares this paper's direction in emphasising difficult examples, but its notion of difficulty is derived from the model and data themselves. The signal used here is external to both: an expert's judgment of the photograph, recorded independently of how any model performs on it. Our fusion module draws on Squeeze-and-Excitation \cite{se} and is used as an established component, not claimed as a contribution.
%%%%%%%%%%%%%%%%%%%%%%%%%%%%%%%%%%%%%%%%%%%%%%%%%%%%%%%%%%%%%%%%
%%%%%%%%%%%%%%%%%%%%%%%%%%%%%%%%%%%%%%%%%%%%%%%%%%%%%%%%%%%%%%%%
\section{Proposed Methodology}
The Proposed architecture is as shown in Fig.~\ref{fig:arch} , it comprises of the following :(i)Image Encoder,(ii)MetaData Encoder (iii)Channel Gated fusion (iv)Prediction Heads (v)Quality-Adaptive Loss  
\begin{figure*}[t]
    \centering
    \includegraphics[width=\textwidth]{mcrsl_archi_widescreen_nolegend.png}
    \vspace{-2mm}
    \caption{Overview of the proposed multimodal architecture. Solid arrows are the forward data path, red dashed arrows the backward gradient flow, and the black dashed arrow the quality path, which enters only the loss and receives no gradient.}
    \label{fig:arch}
\end{figure*}

\subsection{Image Encoder}
\label{sec:image_encoder}
Each dermoscopic image of size $3\times224\times224$ is processed by an EfficientNet-B0 backbone pre-trained on ImageNet \cite{efficientnet,imagenet} and fine-tuned end-to-end. The encoder retains the final convolutional feature map, of size $1280\times7\times7$, rather than only a globally pooled summary; a pooled 1280-dimensional vector can also be derived from this same map by global average pooling. Which representation is used, and at what point pooling occurs, depends on the fusion mechanism, described next.

\subsection{Metadata Encoder}
\label{sec:metadata_encoder}
The metadata associated with each lesion comprises sixteen categorical fields and four numeric fields. Each categorical field is embedded into a twelve-dimensional space, leaving room for a reserved value to handle missing data rather than performing any form of imputation; embedding all sixteen categorical fields results in a 192-dimensional categorical block. The four numeric fields are z-scored based on training-fold statistics, each accompanied by its own binary missingness indicator; the resulting 8-dimensional numeric block is combined with the categorical block to produce a 200-dimensional feature vector, which is fed into a small MLP consisting of a linear layer (200$\to$128), ReLU, dropout with $p=0.2$, a linear layer (128$\to$128), and a final ReLU.

\subsection{Channel-Gated Fusion}
\label{sec:fusion}
Two fusion variants share a common 256-dimensional projection. The late-fusion baseline pools the image feature map directly and concatenates it with the 128-dimensional metadata vector. The channel-gated variant instead uses metadata to modulate the image features before pooling: a linear layer and sigmoid map the metadata representation onto a 1280-dimensional gate,
\begin{equation}
g = \sigma(W_g\, h_{\mathrm{meta}}) \in [0,1]^{1280},
\end{equation}
which rescales the EfficientNet feature map channel-by-channel via elementwise multiplication, $F_{\mathrm{gated}} = g \odot F_{\mathrm{img}}$. Only after this gating does global average pooling followed by a linear/ReLU projection reduce the result to the shared 256-dimensional representation. This is the Squeeze-and-Excitation mechanism \cite{se}, conditioned here on metadata rather than on the feature map itself; it is used as an established component rather than claimed as a contribution of this paper.

\subsection{Prediction Heads}
\label{sec:heads}
Two heads operate on the fused 256-dimensional representation. A binary head, Linear(256$\to$1) followed by a sigmoid,that  predicts $P(\text{malignant})$. A nine-class auxiliary head that predicts the diagnostic subtype and is weighted at 0.4 in the total loss, encouraging the shared representation to retain fine-grained diagnostic structure beyond the binary label.

\subsection{Quality-Adaptive Loss}
\label{sec:qaloss}
The contribution of this work is a per-sample weighting scheme applied to the existing binary
objective. Neither the architecture nor the loss family is modified; what changes is how much
each sample counts within that objective, driven by a real expert-assigned image-quality
rating rather than an estimated quality. Let $q_i \in [1,10]$ denote the mean expert quality
rating assigned to lesion $i$'s diagnostic image. A per-sample multiplicative weight
$w_i^{\mathrm{qual}}$ is applied alongside the existing per-fold class weight
$w_i^{\mathrm{cls}}$:
\begin{equation}
w_i^{\mathrm{cls}} =
\begin{cases}
n_{\mathrm{neg}}/n_{\mathrm{pos}} & y_i = 1 \\
1 & y_i = 0
\end{cases}
\end{equation}
\begin{equation}
\mathcal{L}_i = w_i^{\mathrm{cls}} \cdot w_i^{\mathrm{qual}} \cdot \mathrm{BCE}(\hat{y}_i, y_i).
\end{equation}
Two opposing mappings of the rating onto $[0.5,1.5]$ were compared:
\begin{align}
\text{\emph{trust}:}\quad & w_i^{\mathrm{qual}} = 0.5 + (q_i - 1)/9, \\
\text{\emph{hard\_mining}:}\quad & w_i^{\mathrm{qual}} = 1.5 - (q_i - 1)/9.
\end{align}
\emph{Trust:} It treats low-quality images as less reliable and down-weights them.
\emph{Hard\_mining:} It does the opposite, pushing the model to work harder on degraded images.
Lesions without a valid rating receive a neutral weight of 1.0. $w_i^{\mathrm{qual}}$ is a
precomputed, non-differentiable scalar coefficient, not a learned parameter, so gradients
propagate through the shared architecture by the standard chain rule from the weighted binary
term; only how much each sample contributes changes, never what is differentiated. The weight
range was fixed rather than tuned, and only these two directions were tested, so the
comparison isolates direction from magnitude.

Only \emph{hard\_mining} survived validation, in the opposite direction from AdaFace and MagFace \cite{adaface,magface}. Low-quality samples do not uniformly become less informative under degradation. Image degradation can leave the underlying diagnostic label intact even as the image itself becomes harder to interpret, so training the model to work harder on such a sample does not, in itself, amount to fitting noise. Whether emphasis is warranted depends on what the label was based on and on the nature of the degradation, not on image quality being uniformly detrimental. In dermoscopy, the diagnostic label was assigned from that same image, so the label stays valid and the sample remains an informative hard example rather than a corrupted one.

This mechanism is distinct from the auxiliary quality-regression head described above, which instead asks the model to \emph{predict} the rating as a secondary output, adding parameters and a competing objective to the shared representation. Reweighting adds no parameters and changes only which samples dominate the gradient. The two interventions are mechanistically unrelated, and the auxiliary head's negative result stands independently; it is not superseded by the reweighting result.

Both directions, along with every other configuration compared later, are evaluated under a common subject-disjoint stratified five-fold cross-validation protocol: each subject's lesions are assigned to exactly one fold, so no subject or near-duplicate lesion image appears in both a training and an evaluation fold, and folds are balanced by greedily assigning subjects to minimize the difference in malignant lesion count across folds. For each of the five rotations, one fold is held out as the reported test fold and a second, disjoint fold is held out purely for checkpoint selection by validation balanced accuracy; the test fold never influences model selection, and no hyperparameter is tuned against final test-fold results. This shared protocol is what makes the \emph{trust} versus \emph{hard\_mining} comparison, and the wider ablation grid reported subsequently, directly comparable to one another.
%%%%%%%%%%%%%%%%%%%%%%%%%%%%%%%%%%%%%%%%%%%%%%%%%%%%%%%%%%%%%%%%
\section{Experimental Results}

\subsection{Dataset}
\label{sec:dataset}
MCR-SL \cite{mcrsl} comprises 240 lesions from 60 subjects, 42 malignant and 192 non-malignant, with 1352 dermoscopic and 779 clinical images, 28 lesions histopathology-confirmed, and 17 categorical plus 4 numeric metadata fields. Six lesions carry an undetermined malignancy label and are excluded, leaving 234 usable.
\begin{figure}[htbp]
\centering
\includegraphics[width=\linewidth]{figures/dataset_samples.pdf}
\caption{Representative dermoscopic diagnosis images, one per unified-diagnosis class, with mean expert quality rating $q$.}
\label{fig:samples}
\end{figure}

\subsection{Evaluation Metrics}
The proposed method has been evaluated on the following metrics:
Accuracy is the proportion of all lesions classified correctly, sensitivity the proportion of malignant lesions correctly identified, and specificity the proportion of non-malignant lesions correctly cleared. Since accuracy alone is misleading when the malignant class is a $17.9\%$ minority, the primary metric is balanced accuracy, the unweighted mean of the two class-wise rates, reported alongside macro-averaged F1:
\begin{equation}
\mathrm{BAcc}=\tfrac{1}{2}\left(\mathrm{Sen}+\mathrm{Spe}\right),
\end{equation}
\begin{equation}
\text{macro-F1}=\frac{1}{2}\sum_{c\in\{0,1\}}\frac{2\,\mathrm{Pre}_c\,\mathrm{Rec}_c}{\mathrm{Pre}_c+\mathrm{Rec}_c},
\end{equation}
where $\mathrm{Pre}_c$ and $\mathrm{Rec}_c$ are precision and recall for class $c$. We also report AUROC, the probability that a randomly drawn malignant lesion receives a higher predicted probability than a randomly drawn non-malignant one. All are aggregated across the five test folds as mean and standard deviation. Balanced accuracy has also been used for checkpoint selection.
%%%%%%%%%%%%%%%%%%%%%%%%%%%%%%%%%%%%%%%%%%%%%%%%%%%%%%%%%%%5
\subsection{Results}% REduce text %%
\label{sec:main}
Table~\ref{tab:breakdown} reports the proposed method. Channel-gated fusion has reached the highest macro-F1 but has not dominated: the image-only baseline has higher sensitivity (0.697 vs.\ 0.672) and AUROC (0.895 vs.\ 0.881), and its balanced accuracy is effectively tied (0.784 vs.\ 0.783, against a standard deviation of 0.07--0.09). We report this rather than framing it as a clean win for the fusion method.
\begin{table}[htbp]
\caption{Results of the Proposed Method }
\label{tab:breakdown}
\centering
\scriptsize
\setlength{\tabcolsep}{3pt}
\begin{tabular}{|l|c|c|c|c|}
\hline
\textbf{Step} & \textbf{Quality?} & \textbf{Bal.\ acc.} & \textbf{$\Delta$} & \textbf{Sens.} \\
\hline
Channel-gated baseline & no & 0.783 & --- & 0.672 \\
$+$ \emph{hard\_mining} loss & \textbf{yes} & 0.820 & \textbf{$+$0.037} & 0.764 \\
$+$ test-time augmentation & no & 0.833 & $+$0.050 & \textbf{0.808} \\
$+$ multi-image averaging & no & \textbf{0.840} & $+$0.057 & 0.783 \\
\hline
\emph{hard\_mining}, shuffled & control & 0.775 & $-$0.008 & 0.667 \\
\hline
\end{tabular}
\end{table}

The quality-aware variant, which adds an auxiliary head predicting the quality rating, underperforms on five of six metrics: balanced accuracy 0.740 vs.\ 0.783 and sensitivity 0.578 vs.\ 0.672, against a small specificity gain. Predicting quality is therefore a negative result, mechanistically distinct from reweighting by it.

The proposed quality-adaptive loss improves every core metric over the baseline: balanced accuracy by 0.037 (0.783 to 0.820), sensitivity by 0.092 (0.672 to 0.764), AUROC by 0.025 (0.881 to 0.906) and accuracy to 0.845, against a small specificity cost. The opposite weighting direction is flat on balanced accuracy (0.788) and is not treated as an improvement. The gain is smaller than the fold-to-fold standard deviation of roughly 0.08, and no paired fold-level test reaches significance at $N=5$, so the claim rests on the control below rather than on this comparison.

\smallskip\noindent\textbf{Causal validation via a shuffled-quality control.}
\label{sec:shuffle}
The obvious objection is that any per-sample reweighting might act as a regularizer, leaving the ratings incidental. We have tested this directly: within each fold's training portion the ratings are permuted across lesions, preserving the weight distribution exactly while destroying the lesion-to-rating correspondence. All else is identical, and evaluation always uses the true ratings.

Under this control the proposed loss falls from 0.820 to 0.775 balanced accuracy and from 0.764 to 0.667 sensitivity, returning to the baseline's 0.783 and 0.672. A distribution-matched but uninformative weighting therefore recovers none of the gain, establishing that the improvement depends on the weights tracking real quality rather than on reweighting as such. This is a causal argument rather than a statistical one.

The control has also been informative negatively. The \emph{trust} variant's one apparent advantage, a narrower accuracy gap between the highest and lowest quality terciles (0.082 against the baseline's 0.091), narrows further still under shuffling, to 0.021, while its balanced accuracy falls from 0.788 to 0.763: a weighting carrying no quality information produces a flatter profile than the mechanism intended to produce one. That advantage is therefore reported as a reversed finding rather than a second positive result.
%%%%%%%%%%%%%%%%%%%%%%%%%%%%%%%%%%%%%%%%%%%%%%%%%%%%%%%%%%%%%
\subsection{Qualitative Results} %% To be changed with hard and soft samples and image explanation  % 
\label{sec:qualitative}
Gradients have been taken with respect to the feature map \emph{after} the metadata gate, the representation the decision is computed from, with the metadata branch held fixed. Fig.~\ref{fig:hardsamples} contrasts two regimes chosen by a fixed rule rather than by inspection: the top row holds the three lesions both models call most confidently, the bottom row the three lowest-rated lesions the baseline gets wrong and the proposed loss corrects.

The rows separate cleanly by quality, mean expert rating 7.1 against 4.4, with all three corrected lesions in the lowest tercile where the loss assigns its largest weights. In the top row attention sits compactly on the lesion and the two models agree to within $10^{-3}$, so reweighting leaves easy cases untouched. In the bottom row it changes the call: for the melanoma in (e), rated 4.3, baseline attention falls away from the pigmented region and the lesion is missed at $P=0.28$, whereas under reweighting attention moves onto the lesion and the call becomes $P=0.89$. The two false positives in (d) and (f) are also corrected, though (d) only marginally, from $P=0.53$ to $P=0.50$.

These are illustrative rather than representative. Across the 231 labelled lesions the reweighting fixes six malignant lesions and breaks two, a net gain of four, which is what a $+0.092$ sensitivity amounts to at this scale; on non-malignant lesions it is net negative, seven fixed against twelve broken, which is the specificity cost noted above.

\begin{figure}[htbp]
\centering
\includegraphics[width=\linewidth]{figures/hard_vs_nonhard_gradcam.pdf}
\caption{Grad-CAM on the post-gate feature map. Top: lesions both models call confidently. Bottom: low-rated lesions the baseline misclassifies and the proposed loss corrects.}
\label{fig:hardsamples}
\end{figure}


\subsection{Comparison with the Literature}
\label{sec:comparison}
MCR-SL has had no prior published benchmark, so the results reported here are the first on this dataset. This is not framed as a state-of-the-art claim, since there is no prior number to have surpassed; the loss mechanism and its validation remain the principal contribution, and the comparison below is offered only to situate that result against the closest available datasets and published binary-task work. Table~\ref{tab:literature} summarizes this comparison.

\begin{table}[htbp]
\caption{Comparison with existing works.}
\label{tab:literature}
\centering
\resizebox{\columnwidth}{!}{%
\begin{tabular}{|l|l|c|c|c|c|}
\hline
\textbf{Method} & \textbf{Data} & \textbf{N} & \textbf{Mal.\%} & \textbf{Task} & \textbf{Bal.acc} \\
\hline
Baseline (ours) & MCR-SL & 234 & 17.9 & bin. & 0.783 \\
$+$hard\_mining (ours) & MCR-SL & 234 & 17.9 & bin. & 0.820 \\
Best config (ours) & MCR-SL & 234 & 17.9 & bin. & \textbf{0.840} \\
\hline
Uliana \& Krohling \cite{diffmic} & PAD-UFES-20 & 2298 & 47.5 & bin. & 0.836 \\
\hline
PAD-UFES-20 (native) & PAD-UFES-20 & --- & --- & 6-cls & 0.77--0.84 \\
HAM10000 \cite{ham10000} & HAM10000 & $\sim$40$\times$ & --- & derm. & high acc.\\
\hline
\end{tabular}%
}
\end{table}
The only directly comparable binary-task result is Uliana and Krohling's \cite{diffmic} (0.836 balanced accuracy, diffusion classifier, 2298 images, five-fold cross-validation); the 0.004-point gap to this paper's best configuration (0.840) is far smaller than the fold-to-fold standard deviation, so the two results are statistically indistinguishable, and the datasets also differ substantially in unit (lesions vs.\ images) and class balance (17.9\% vs.\ 47.5\% malignant). PAD-UFES-20's native six-class balanced accuracies (0.765--0.842) and HAM10000's \cite{ham10000} higher raw accuracy are shown for context only and are not valid comparators, since neither shares this paper's binary task, metadata use, or dataset scale.
\subsection{Ablation and Robustness Analyses}
\begin{table}[htbp]
\caption{Ablation Results of the Proposed Fusion Mechanism  }
\label{tab:ablation}
\centering
\scriptsize
\setlength{\tabcolsep}{2pt}
\begin{tabular}{|l|c|c|c|}
\hline
\textbf{Metric} & \textbf{Image-only} & \textbf{Late fusion} & \textbf{Channel-gated} \\
\hline
Accuracy & 0.813 $\pm$ 0.020 & 0.831 $\pm$ 0.052 & \textbf{0.833 $\pm$ 0.070} \\
Balanced acc. & \textbf{0.784 $\pm$ 0.073} & 0.768 $\pm$ 0.090 & 0.783 $\pm$ 0.085 \\
Macro F1 & 0.733 $\pm$ 0.038 & 0.743 $\pm$ 0.051 & \textbf{0.757 $\pm$ 0.073} \\
Sensitivity & \textbf{0.697 $\pm$ 0.187} & 0.644 $\pm$ 0.191 & 0.672 $\pm$ 0.159 \\
Specificity & 0.872 $\pm$ 0.061 & 0.891 $\pm$ 0.021 & \textbf{0.895 $\pm$ 0.048} \\
AUROC & \textbf{0.895 $\pm$ 0.053} & 0.851 $\pm$ 0.080 & 0.881 $\pm$ 0.065 \\
\hline
\end{tabular}
\end{table}

Eleven configurations have been run under an identical protocol, each varying one component of
the channel-gated method: the loss (focal loss \cite{focal}), the input (hair removal and colour
normalisation), the objective (a contrastive auxiliary term \cite{supcon}), the optimiser
(decoupled weight decay with a cosine schedule \cite{adamw}, sharpness-aware minimisation
\cite{sam}), and the evaluation (flip-averaged and multi-image test-time averaging), plus four
stacked combinations. Table~\ref{tab:sweep} reports all eleven. Sharpness-aware minimisation with
test-time augmentation is the best of the sweep at 0.822; the contrastive term and the alternative
optimiser each underperform the baseline, and stacking every promising intervention reaches only
0.759.

Applying the two evaluation-time interventions on top of the proposed loss instead has produced
this paper's best configuration, decomposed in Table~\ref{tab:breakdown}: test-time augmentation
has raised balanced accuracy to 0.833 and sensitivity to 0.808, and multi-image averaging has
raised it further to 0.840 $\pm$ 0.095, stacking cleanly here unlike in the sweep above. Of the
total 0.057 gain, 0.037, roughly 65\%, has come from the quality-adaptive loss and is causally
validated above; the remaining 0.020 has come from generic evaluation-time averaging that needs
no retraining and has nothing to do with image quality. We state this split so that the
mechanism's contribution is not overstated by the final figure.

\begin{table}[htbp]
\caption{Ablation Results on Hyper parameters }
\label{tab:sweep}
\centering
\tiny
\setlength{\tabcolsep}{3pt}
\begin{tabular}{|l|c|c|c|c|c|}
\hline
\textbf{Configuration} & \textbf{Acc.} & \textbf{Bal.\ acc.} & \textbf{Sens.} & \textbf{Spec.} & \textbf{AUROC} \\
\hline
Focal loss & 0.858 & 0.809 & 0.694 & 0.923 & 0.877 \\
Dermoscopy preprocessing & 0.828 & 0.793 & \textbf{0.736} & 0.851 & 0.875 \\
Contrastive auxiliary loss & 0.812 & 0.770 & 0.669 & 0.872 & 0.873 \\
Alt.\ optimizer (AdamW+cosine) & 0.822 & 0.757 & 0.661 & 0.853 & 0.845 \\
Test-time augmentation (TTA) & 0.841 & 0.788 & 0.672 & 0.904 & 0.897 \\
Multi-image averaging & 0.854 & 0.813 & 0.719 & 0.907 & 0.906 \\
Sharpness-aware min.\ (SAM) & 0.858 & 0.810 & 0.694 & 0.927 & 0.904 \\
SAM + multi-image & 0.851 & 0.800 & 0.672 & 0.928 & \textbf{0.917} \\
\textbf{SAM + TTA} & \textbf{0.861} & \textbf{0.822} & 0.719 & 0.925 & 0.908 \\
SAM + TTA + multi-image & 0.856 & 0.810 & 0.694 & 0.925 & 0.915 \\
Preprocess.+focal+SAM+TTA & 0.818 & 0.759 & 0.622 & 0.895 & 0.893 \\
\hline
\end{tabular}
\end{table}




\label{sec:robustness}
These analyses use the channel-gated out-of-fold predictions over the 231 lesions carrying both a prediction and a quality rating; three of the 234 lacked a usable image in their assigned test fold and were excluded rather than imputed.

\textbf{Quality-stratified performance.} Terciles by mean expert rating give accuracies of 0.776, 0.914 and 0.868 (low, mid, high; sizes 85, 93, 53, uneven because the mean of up to three integer ratings takes few distinct values). The pattern is not monotonic and Spearman correlation between rating and error is $-0.093$ ($p=0.158$): no strong quality-performance relationship is detected at 231 lesions, neither confirming nor ruling one out.

The validated quality-adaptive loss has not flattened this profile: its high-minus-low accuracy gap has remained 0.091. Its benefit has appeared within terciles instead, most visibly on the lowest-quality one, where malignant sensitivity has risen from 0.667 to 0.810. The mechanism has improved detection on poor images without equalizing accuracy across quality levels, and we have not claimed the latter.

\textbf{Reversed findings.} Three apparently positive results have not survived scrutiny and are reported rather than omitted. \emph{Trust}'s narrower quality-tercile gap was undercut by the shuffled control narrowing it further, as noted earlier in the results. A pooled correlation between expert certainty and prediction error, initially $\rho=-0.175$ ($p=0.008$, $N=231$), dissolved once stratified by class ($p=0.083$ and $p=0.559$), having reflected class composition rather than an independent relationship. And a proposed change to train on all images per lesion proved on verification to describe what the pipeline already did. At this sample size an uncontrolled first look repeatedly produces positives that controls dissolve, and reporting that pattern is part of the contribution.

\textbf{Intra-subject consistency.} Of 55 subjects with two or more evaluable lesions, 27 were perfectly consistent against a permutation null mean of 29.9 ($p=0.954$), and none had all lesions misclassified: errors have not concentrated in particular patients.

\textbf{Histopathology-confirmed versus panel-consensus.} The model has performed worse on the 28 histopathology-confirmed lesions than on the 203 panel-consensus-only ones, 0.750 against 0.867 accuracy, plausibly because histopathology is more often ordered for clinically ambiguous lesions. With 28 lesions this has been treated as qualitative only.



\section{Conclusion and Future Work}
Expert quality judgments are rarely recorded alongside diagnostic labels, so quality-aware training has typically relied on estimated proxies rather than direct observation. This paper proposes quality-adaptive loss reweighting, scaling each sample's loss by a real expert quality rating, and finds that up-weighting low-quality samples — contrary to face-recognition precedent — improves balanced accuracy by 0.037 and sensitivity by 0.092 over baseline. As this gain falls within fold-to-fold variance, it was validated causally: a shuffled-quality control removed the gain entirely, confirming the improvement depends on real quality information rather than reweighting alone. The full pipeline reaches 0.840 balanced accuracy, roughly 65\% attributable to the loss mechanism. This work was enabled by MCR-SL, the only public skin lesion dataset with per-image expert quality ratings, on which a first benchmark and six robustness analyses are also reported.



%%%%%%%%%%%%%%%%%%%%%%%%%%%%%%%%%%%%%%%%%%%%%%%%%%%%%%%%%%%%%%
\begin{thebibliography}{00}
\bibitem{mcrsl} M. Castro-Fernandez \emph{et al.}, ``MCR-SL: A multimodal, context-rich skin lesion dataset for skin cancer diagnosis,'' \emph{Data}, vol. 10, no. 10, art. 166, 2025.
\bibitem{esteva2017} A. Esteva, B. Kuprel, R. A. Novoa, J. Ko, S. M. Swetter, H. M. Blau, and S. Thrun, ``Dermatologist-level classification of skin cancer with deep neural networks,'' \emph{Nature}, vol. 542, pp. 115--118, 2017.
\bibitem{ham10000} P. Tschandl, C. Rosendahl, and H. Kittler, ``The HAM10000 dataset, a large collection of multi-source dermatoscopic images of common pigmented skin lesions,'' \emph{Sci. Data}, vol. 5, art. 180161, 2018.
\bibitem{metablock} A. G. C. Pacheco and R. A. Krohling, ``An attention-based mechanism to combine images and metadata in deep learning models applied to skin cancer classification,'' \emph{IEEE J. Biomed. Health Inform.}, vol. 25, no. 9, pp. 3554--3563, 2021.
\bibitem{paduf20} A. G. C. Pacheco \emph{et al.}, ``PAD-UFES-20: A skin lesion dataset composed of patient data and clinical images collected from smartphones,'' \emph{Data in Brief}, vol. 32, art. 106221, 2020.
\bibitem{se} J. Hu, L. Shen, and G. Sun, ``Squeeze-and-excitation networks,'' in \emph{Proc. IEEE Conf. Computer Vision and Pattern Recognition (CVPR)}, 2018.
\bibitem{efficientnet} M. Tan and Q. V. Le, ``EfficientNet: Rethinking model scaling for convolutional neural networks,'' in \emph{Proc. Int. Conf. Machine Learning (ICML)}, 2019, pp. 6105--6114.
\bibitem{imagenet} J. Deng, W. Dong, R. Socher, L.-J. Li, K. Li, and L. Fei-Fei, ``ImageNet: A large-scale hierarchical image database,'' in \emph{Proc. IEEE Conf. Computer Vision and Pattern Recognition (CVPR)}, 2009, pp. 248--255.
\bibitem{sam} P. Foret, A. Kleiner, H. Mobahi, and B. Neyshabur, ``Sharpness-aware minimization for efficiently improving generalization,'' in \emph{Proc. Int. Conf. Learning Representations (ICLR)}, 2021.
\bibitem{adamw} I. Loshchilov and F. Hutter, ``Decoupled weight decay regularization,'' in \emph{Proc. Int. Conf. Learning Representations (ICLR)}, 2019.
\bibitem{supcon} P. Khosla \emph{et al.}, ``Supervised contrastive learning,'' in \emph{Advances in Neural Information Processing Systems (NeurIPS)}, vol. 33, 2020, pp. 18661--18673.
\bibitem{focal} T.-Y. Lin, P. Goyal, R. Girshick, K. He, and P. Doll\'ar, ``Focal loss for dense object detection,'' in \emph{Proc. IEEE Int. Conf. Computer Vision (ICCV)}, 2017, pp. 2980--2988.
\bibitem{diffmic} J. J. M. Uliana and R. A. Krohling, ``Diffusion models applied to skin and oral cancer classification,'' 2025, arXiv:2504.00026. [Online]. Available: https://arxiv.org/abs/2504.00026
\bibitem{adaface} M. Kim, A. K. Jain, and X. Liu, ``AdaFace: Quality adaptive margin for face recognition,'' in \emph{Proc. IEEE/CVF Conf. Computer Vision and Pattern Recognition (CVPR)}, 2022.
\bibitem{magface} Q. Meng, S. Zhao, Z. Huang, and F. Zhou, ``MagFace: A universal representation for face recognition and quality assessment,'' in \emph{Proc. IEEE/CVF Conf. Computer Vision and Pattern Recognition (CVPR)}, 2021.
\bibitem{hu2023} M. Hu and X. Yang, ``Attention-driven lightweight model for pigmented skin lesion detection,'' 2023, arXiv:2308.02119. [Online]. Available: https://arxiv.org/abs/2308.02119
\bibitem{qgdr} Y. Pan, X. Cai, P. Xiong, F. Mei, Z. Wang, and J. Hong, ``From discarding to leveraging: quality-aware collaborative learning for robust diabetic retinopathy grading,'' \emph{Frontiers in Medicine}, vol. 13, art. 1795484, 2026.
\bibitem{deepdrid} R. Liu, X. Wang, Q. Wu \emph{et al.}, ``DeepDRiD: Diabetic retinopathy---grading and image quality estimation challenge,'' \emph{Patterns}, vol. 3, no. 6, art. 100512, 2022.
\end{thebibliography}

\end{document}
